% !TEX root = ../../../main.tex
\section{Kernel}
\label{sec:linear-algebra-i:kernel}

% Sources: supplied lecture photos and September 24 slides 3 and 12.
\begin{definition}
  Let \(V\) and \(W\) be vector spaces and \(L:V\to W\) a linear map.
  Its kernel is\marginnote{%
    \emph{Kernel} comes from Old English \emph{cyrnel}, meaning a seed or
    pip, related to \emph{corn} with a diminutive ending. The word later
    acquired the figurative sense of a central core \cite{EtymologyKernel}.
    Its technical meaning depends on the context.}
  \[
    \ker L=\{v\in V:L(v)=0_W\} \subseteq V.
  \]
  The kernel is a subspace of the domain \(V\), consisting of the inputs
  mapped to the zero vector.
\end{definition}

Figure~\ref{fig:linear-algebra-i:kernel-to-zero} emphasizes the defining
property of the kernel. Every input in this subspace has the same output,
the zero vector \(0_W\).

\begin{figure}[htbp]
  \centering
  % Kernel-only adaptation of Pillsmarch's kernel/image illustration.
% Source: https://commons.wikimedia.org/wiki/File:Kernel_and_image_of_linear_map.svg
% Original: Pillsmarch, December 20, 2023. License: CC BY 4.0.
% This schematic omits the image and the surrounding domain/codomain.
% Diagram adaptation licensed under https://creativecommons.org/licenses/by/4.0/
\begin{tikzpicture}[
  x=1mm,
  y=1mm,
  font=\footnotesize,
  line width=0.45pt,
  every node/.style={inner sep=2pt},
  map arrow/.style={line width=0.35pt,-{Latex[length=1.6mm,width=1.1mm]}}
]
  \filldraw[fill=black!15] (15,0) circle[radius=13mm];
  \node at (15,0) {\(\ker L\)};
  \node at (46,18) {\(L:V\to W\)};

  \coordinate (codomain-zero) at (73,0);
  \draw[map arrow] (21.5,11.2583)
    .. controls (39,11) and (56,5) .. (codomain-zero);
  \draw[map arrow] (28,0) -- (codomain-zero);
  \draw[map arrow] (21.5,-11.2583)
    .. controls (39,-11) and (56,-5) .. (codomain-zero);

  \fill (codomain-zero) circle[radius=1.6pt];
  \node at (78,3) {\(0_W\)};
\end{tikzpicture}

  \caption{The kernel of \(L:V\to W\), shown schematically.
    All vectors in the shaded region map to \(0_W\).
    The circle indicates a set, not a bounded vector space.
    TikZ redrawing inspired by
    Wikipedia's \emph{Kernel (linear algebra)} page
    \cite{WikipediaKernel} and Pillsmarch's illustration
    \cite{PillsmarchKernelImage},
    licensed under \href{https://creativecommons.org/licenses/by/4.0/}{CC BY 4.0}.}
  \label{fig:linear-algebra-i:kernel-to-zero}
\end{figure}

\subsection{Null space, nullity, and homogeneous equations}

The kernel is also called the \emph{null space}. Here ``null'' refers to
the zero output, while ``space'' refers to the subspace of inputs that
produce it. A vector in the null space need not itself be zero.

The \emph{nullity} of \(L\) is the dimension of its null space,
\[
  \operatorname{nullity}L=\dim\ker L.
\]
When the kernel is finite-dimensional, its nullity is the number of vectors
in any basis of \(\ker L\). The null space is a subspace of inputs;
nullity counts its independent directions.

% The author requested explicit rank/nullity bounds on October 1, 2026.
In finite dimensions, nullity is therefore a nonnegative integer.
If \(\ker L=\{0_V\}\), its basis is empty, so its nullity is \(0\).
The zero vector is still in the kernel, but it contributes no independent
direction. Since the kernel is a subspace of \(V\), for finite-dimensional
\(V\) we have
\[
  0\leq\operatorname{nullity}L\leq\dim V.
\]

For a matrix \(A\in\mathbb F^{n\times m}\), the same definition becomes
\[
  \operatorname{Null}(A)=\ker A
    =\{x\in\mathbb F^m:Ax=0\}.
\]
Thus finding the null space means solving a homogeneous system of linear
equations \cite{MargalitRabinoffSubspaces}, using the row operations of
Section~\ref{sec:linear-algebra-i:systems-of-linear-equations}.
The dimension of this solution space is \(\operatorname{nullity}A\).

\subsection{Information lost by the map}

The kernel records changes to an input that leave its output unchanged.
If \(k\in\ker L\), then
\[
  L(v+k)=L(v)+L(k)=L(v).
\]
Conversely, if \(L(v_1)=L(v_2)\), linearity gives
\[
  L(v_1-v_2)=L(v_1)-L(v_2)=0_W,
\]
so \(v_1-v_2\in\ker L\). Two inputs are indistinguishable under \(L\)
exactly when their difference lies in its kernel.

This matters whenever a linear map models measurements. An input can change
by any vector in the kernel without changing the observed data.
\marginnote{%
  In computer vision, a \emph{kernel} usually means an array of weights
  used for convolution or correlation \cite{OpenCVLinearFilters}.
  The array defines a filter. For a linear filter, its linear-algebra
  kernel would instead be the subspace of images it maps to zero.}
Nullity therefore counts the independent directions of this ambiguity.
Section~\ref{sec:linear-algebra-i:rank} defines rank as the dimension of
the image. For finite-dimensional \(V\), the rank--nullity theorem
(Theorem~\ref{thm:linear-algebra-i:rank-nullity}) gives
\[
  \operatorname{rank}L+\operatorname{nullity}L=\dim V.
\]
The independent output directions, counted by rank, and the input
directions lost in the kernel, counted by nullity, together account for
the dimension of the domain.

More generally, suppose \(v_0\) is one solution of \(L(v)=w\). Then all
solutions are precisely
\[
  v_0+\ker L=\{v_0+k:k\in\ker L\}.
\]
Indeed, adding \(k\in\ker L\) preserves the output, and any other solution
differs from \(v_0\) by a kernel vector. For a matrix equation \(Ax=b\),
this says that every solution is a particular solution plus a solution
of \(Ax=0\) \cite{MargalitRabinoffSolutionSets}. The image determines
whether a solution exists; the kernel determines whether an existing
solution is unique.

\subsection{Injectivity}

The general definition of injectivity is recalled in
\S~\ref{sec:appendix:map-properties}.
The map \(L\) is injective exactly when \(\ker L=\{0_V\}\), or,
in finite dimensions, when \(\dim\ker L=0\). If \(L\) is injective,
\(L(k)=0_W=L(0_V)\) forces \(k=0_V\) for every \(k\in\ker L\).
Conversely, for a zero kernel, injectivity follows from the equivalences
\[
  \begin{aligned}
    L(v_1)=L(v_2)
      &\iff L(v_2)-L(v_1)=0_W\\
      &\iff L(v_2-v_1)=0_W\\
      &\iff v_2-v_1\in\ker L=\{0_V\}\\
      &\iff v_2=v_1.
  \end{aligned}
\]
% The zero-kernel chain follows the supplied material. The converse and
% null-space context were added at the author's request on September 30, 2026.
